\paragrafos{Ορθογώνια, ρόμβοι, τετράγωνα}{Β.3.11, «Είδη παραλληλογράμμων και ιδιότητες: ορθογώνια, ρόμβοι, τετράγωνα»}

\begin{anagke}
Η Αγγελική διπλώνει χαρτί σε τρία σχήματα: ένα ορθογώνιο, έναν ρόμβο και ένα τετράγωνο. Για το καθένα ψάχνει με πόσους τρόπους μπορεί να το διπλώσει ώστε τα δύο μισά να κάθονται ακριβώς το ένα πάνω στο άλλο. «Το ορθογώνιο δύο, ο ρόμβος δύο, το τετράγωνο τέσσερις!», λέει. Ο Στέλιος αναφωνεί: «Το τετράγωνο τα έχει όλα, άρα ό,τι ισχύει για τα άλλα ισχύει και για αυτό!». Η Ρήνη όμως ρωτά: «Δηλαδή το τετράγωνο είναι και ορθογώνιο και ρόμβος; Δεν βγάζει νόημα». Η Φανή γράφει στο τετράδιό της και ρωτά τι ακριβώς κάνει κάθε σχήμα «ξεχωριστό».

Σε αυτή την παράγραφο ορίζουμε το ορθογώνιο, τον ρόμβο και το τετράγωνο, βρίσκουμε τις ιδιότητές τους και καταλαβαίνουμε πώς σχετίζονται.
\end{anagke}

\begin{lexi}{ορθογώνιο, ρόμβος, τετράγωνο}{ονόματα παραλληλογράμμων}
Το \lexiword{ορθογώνιο} έχει όλες τις γωνίες \emph{ορθές} (από το «ορθός» και «γωνία»). Ο \lexiword{ρόμβος} ήταν στην αρχαιότητα η σβούρα, που «γυρίζει». Το \lexiword{τετράγωνο} έχει «τέσσερις γωνίες» ίσες, και τέσσερις ίσες πλευρές.
\end{lexi}

\begin{metafrasi}
Παραλληλόγραμμο με όλες τις γωνίες ορθές. & ορθογώνιο\\
Παραλληλόγραμμο με όλες τις πλευρές ίσες. & ρόμβος\\
Ορθογώνιο και ρόμβος μαζί. & τετράγωνο\\
Η ευθεία που διπλώνει το σχήμα σε δύο ίσα μέρη. & άξονας συμμετρίας\\
\end{metafrasi}

\tmima{Η θεωρία}

Τα τρία αυτά σχήματα είναι παραλληλόγραμμα, άρα έχουν όλες τις ιδιότητες του παραλληλογράμμου (ίσες απέναντι πλευρές και γωνίες, διχοτομούμενες διαγώνιες). Έχουν όμως και κάτι παραπάνω.

\begin{orismos}{ορθογώνιο}
\textbf{Ορθογώνιο} (παραλληλόγραμμο) λέγεται το παραλληλόγραμμο που έχει όλες τις γωνίες του ορθές.
\end{orismos}

\begin{idiotita}{ορθογωνίου}
Εκτός από τις ιδιότητες του παραλληλογράμμου, στο ορθογώνιο:
\begin{itemize}[leftmargin=1.2em,itemsep=0.3ex,topsep=2pt]
\item οι \textbf{διαγώνιοι είναι ίσες}: $\mathrm{ΑΓ}=\mathrm{ΒΔ}$,
\item οι \textbf{μεσοκάθετοι των πλευρών} του είναι \textbf{άξονες συμμετρίας} (έχει δύο άξονες).
\end{itemize}
\end{idiotita}

\begin{orismos}{ρόμβος}
\textbf{Ρόμβος} λέγεται το παραλληλόγραμμο που έχει όλες τις πλευρές του ίσες.
\end{orismos}

\begin{idiotita}{ρόμβου}
Εκτός από τις ιδιότητες του παραλληλογράμμου, στον ρόμβο:
\begin{itemize}[leftmargin=1.2em,itemsep=0.3ex,topsep=2pt]
\item οι \textbf{διαγώνιοι τέμνονται κάθετα}: $\mathrm{ΑΓ}\perp\mathrm{ΒΔ}$,
\item οι \textbf{διαγώνιοι διχοτομούν τις γωνίες} του,
\item οι \textbf{ευθείες των διαγωνίων} είναι \textbf{άξονες συμμετρίας} (έχει δύο άξονες).
\end{itemize}
\end{idiotita}

\begin{orismos}{τετράγωνο}
\textbf{Τετράγωνο} λέγεται το παραλληλόγραμμο που έχει όλες τις γωνίες του ορθές και όλες τις πλευρές του ίσες.
\end{orismos}

\begin{idiotita}{τετραγώνου}
Το τετράγωνο είναι ταυτόχρονα \textbf{ορθογώνιο και ρόμβος}, άρα έχει όλες τις ιδιότητές τους. Οι διαγώνιοι του είναι ίσες, κάθετες και διχοτομούν τις γωνίες του (κάθε γωνία της διαγωνίου με την πλευρά είναι $45^\circ$). Έχει \textbf{τέσσερις άξονες συμμετρίας}: τις δύο διαγωνίους και τις δύο μεσοκαθέτους των πλευρών.
\end{idiotita}

\begin{center}
\begin{tikzpicture}[scale=0.62]
\draw[thick] (0,0)--(4,0)--(4,2.6)--(0,2.6)--cycle;
\draw[gblue,dashed] (-0.4,1.3)--(4.4,1.3); \draw[gblue,dashed] (2,-0.4)--(2,3);
\node at (2,-1) {\small ορθογώνιο};
\begin{scope}[xshift=7cm]
\draw[thick] (2,-0.3)--(4,1.3)--(2,2.9)--(0,1.3)--cycle;
\draw[gblue,dashed] (-0.4,1.3)--(4.4,1.3); \draw[gblue,dashed] (2,-0.7)--(2,3.3);
\node at (2,-1) {\small ρόμβος};
\end{scope}
\begin{scope}[xshift=14cm]
\draw[thick] (0.2,0)--(3.8,0)--(3.8,3.6)--(0.2,3.6)--cycle;
\draw[gblue,dashed] (0,0)--(4,3.6); \draw[gblue,dashed] (0,3.6)--(4,0);
\draw[gblue,dashed] (-0.2,1.8)--(4.2,1.8); \draw[gblue,dashed] (2,-0.2)--(2,3.8);
\node at (2,-1) {\small τετράγωνο};
\end{scope}
\end{tikzpicture}
\end{center}

Το διάγραμμα δείχνει πώς συνδέονται όλα τα τετράπλευρα που μάθαμε.

\begin{center}
\begin{tikzpicture}[scale=0.9,node distance=0.8cm,
  box/.style={draw,rounded corners=2pt,fill=gblueL,inner sep=4pt,font=\small,align=center}]
\node[box] (t) at (3,3.6) {τετράπλευρα};
\node[box] (p) at (0.8,2.1) {παραλληλόγραμμα};
\node[box] (tr) at (5.2,2.1) {τραπέζια};
\node[box] (o) at (-0.6,0.6) {ορθογώνια};
\node[box] (r) at (3,0.6) {ρόμβοι};
\node[box] (q) at (1.2,-0.8) {τετράγωνα};
\draw[->,thick] (t)--(p); \draw[->,thick] (t)--(tr);
\draw[->,thick] (p)--(o); \draw[->,thick] (p)--(r);
\draw[->,thick] (o)--(q); \draw[->,thick] (r)--(q);
\end{tikzpicture}
\end{center}

Οι διαγώνιοι χωρίζουν τα σχήματα σε τρίγωνα: στο ορθογώνιο σε τέσσερα ισοσκελή τρίγωνα (γιατί $\mathrm{ΟΑ}=\mathrm{ΟΒ}=\mathrm{ΟΓ}=\mathrm{ΟΔ}$), στον ρόμβο σε τέσσερα ορθογώνια τρίγωνα (γιατί οι διαγώνιοι είναι κάθετες) και στο τετράγωνο σε τέσσερα ορθογώνια και ισοσκελή τρίγωνα. Η περίμετρος του ορθογωνίου είναι $2\cdot(\mathrm{ΑΒ}+\mathrm{ΒΓ})$ και του ρόμβου και του τετραγώνου $4\cdot(\text{πλευρά})$.

Σχετικά με το ύψος: όλες οι πλευρές του ορθογωνίου και του τετραγώνου είναι και ύψη, γιατί είναι κάθετες μεταξύ τους. Στον ρόμβο τα ύψη που φέρνουμε από μια κορυφή είναι ίσα, ενώ σε ένα τυχαίο παραλληλόγραμμο δεν είναι αναγκαστικά ίσα.

\begin{prosoxi}
\begin{itemize}[leftmargin=1.2em,itemsep=0.3ex,topsep=2pt]
\item Κάθε τετράγωνο είναι ρόμβος και ορθογώνιο, αλλά \emph{όχι} κάθε ρόμβος ή ορθογώνιο είναι τετράγωνο.
\item Οι διαγώνιοι του ορθογωνίου είναι ίσες αλλά \emph{όχι κάθετες}, ενώ του ρόμβου είναι κάθετες αλλά \emph{όχι ίσες}.
\item Το ορθογώνιο έχει δύο άξονες συμμετρίας και όχι τέσσερις· οι διαγώνιοί του \emph{δεν} είναι άξονες.
\item Ο ρόμβος δεν έχει ορθές γωνίες, εκτός αν είναι τετράγωνο.
\end{itemize}
\end{prosoxi}

\tmima{Λυμένα παραδείγματα}

\paradeigma{Το ορθογώνιο}
\tixero Ορθογώνιο $\mathrm{ΑΒΓΔ}$ με $\mathrm{ΑΒ}=8$ cm, $\mathrm{ΒΓ}=6$ cm και $\mathrm{ΑΓ}=10$ cm.\par
\tithelo Η περίμετρος, η $\mathrm{ΒΔ}$ και τα τμήματα $\mathrm{ΟΑ}$, $\mathrm{ΟΒ}$ ($\mathrm{Ο}$ το κέντρο).\par
\possindeo Η περίμετρος είναι $2\cdot(8+6)=28$ cm. Οι διαγώνιοι του ορθογωνίου είναι ίσες, άρα $\mathrm{ΒΔ}=\mathrm{ΑΓ}=10$ cm. Διχοτομούνται, άρα $\mathrm{ΟΑ}=\mathrm{ΟΒ}=10:2=5$ cm.\par
\apantisi Περίμετρος $28$ cm, $\mathrm{ΒΔ}=10$ cm, $\mathrm{ΟΑ}=\mathrm{ΟΒ}=5$ cm.\par
\elegxos Το τρίγωνο $\mathrm{ΟΑΒ}$ είναι ισοσκελές ($\mathrm{ΟΑ}=\mathrm{ΟΒ}$).

\paradeigma{Ο ρόμβος και οι γωνίες του}
\tixero Ρόμβος $\mathrm{ΑΒΓΔ}$ με $\hat{\mathrm{Α}}=40^\circ$.\par
\tithelo Οι γωνίες του και η γωνία $\widehat{\mathrm{ΑΒΟ}}$, όπου $\mathrm{Ο}$ το σημείο τομής των διαγωνίων.\par
\possindeo Η $\hat{\mathrm{Γ}}=40^\circ$ (απέναντι) και $\hat{\mathrm{Β}}=\hat{\mathrm{Δ}}=180^\circ-40^\circ=140^\circ$. Η διαγώνιος $\mathrm{ΑΓ}$ διχοτομεί την $\hat{\mathrm{Α}}$, άρα $\widehat{\mathrm{ΟΑΒ}}=20^\circ$. Οι διαγώνιοι είναι κάθετες, άρα $\widehat{\mathrm{ΑΟΒ}}=90^\circ$. Στο τρίγωνο $\mathrm{ΑΟΒ}$: $\widehat{\mathrm{ΑΒΟ}}=180^\circ-20^\circ-90^\circ=70^\circ$.\par
\apantisi $40^\circ,140^\circ,40^\circ,140^\circ$ και $\widehat{\mathrm{ΑΒΟ}}=70^\circ$.\par
\elegxos Η $\widehat{\mathrm{ΑΒΟ}}$ είναι η μισή της $\hat{\mathrm{Β}}$: $140^\circ:2=70^\circ$.

\paradeigma{Ρόμβος με διαγωνίους}
\tixero Ρόμβος με πλευρά $6{,}5$ cm και διαγωνίους $12$ cm και $5$ cm.\par
\tithelo Η περίμετρος του ρόμβου και η περίμετρος του τριγώνου $\mathrm{ΑΟΒ}$.\par
\possindeo Η περίμετρος του ρόμβου είναι $4\cdot6{,}5=26$ cm. Οι διαγώνιοι διχοτομούνται: $\mathrm{ΟΑ}=12:2=6$ cm και $\mathrm{ΟΒ}=5:2=2{,}5$ cm. Το τρίγωνο $\mathrm{ΑΟΒ}$ έχει πλευρές $6$, $2{,}5$ και $6{,}5$.\par
\apantisi $26$ cm και $6+2{,}5+6{,}5=15$ cm.\par
\elegxos Η περίμετρος του ρόμβου είναι $4$ φορές την πλευρά.

\paradeigma{Το τετράγωνο με άγνωστο}
\tixero Στο τετράγωνο $\mathrm{ΑΒΓΔ}$ είναι $\mathrm{ΑΒ}=3x+1$ και $\mathrm{ΒΓ}=2x+6$.\par
\tithelo Το $x$, η πλευρά και η περίμετρος.\par
\possindeo Όλες οι πλευρές του τετραγώνου είναι ίσες, άρα $3x+1=2x+6$, δηλαδή $x=5$. Η πλευρά είναι $3\cdot5+1=16$ και η περίμετρος $4\cdot16=64$.\par
\apantisi $x=5$, πλευρά $16$, περίμετρος $64$.\par
\elegxos $2\cdot5+6=16$.

\paradeigma{Τρίγωνα από τις διαγωνίους}
\tixero Ένα ορθογώνιο, ένας ρόμβος και ένα τετράγωνο με τις διαγωνίους τους.\par
\tithelo Σε τι τρίγωνα χωρίζονται;\par
\possindeo Στο ορθογώνιο οι διαγώνιοι είναι ίσες και διχοτομούνται, άρα τα τέσσερα τρίγωνα έχουν δύο ίσες πλευρές: ισοσκελή. Στον ρόμβο οι διαγώνιοι είναι κάθετες: ορθογώνια τρίγωνα. Στο τετράγωνο και τα δύο: ορθογώνια και ισοσκελή.\par
\apantisi Ορθογώνιο: ισοσκελή. Ρόμβος: ορθογώνια. Τετράγωνο: ορθογώνια και ισοσκελή.\par
\elegxos Το τετράγωνο έχει τις ιδιότητες και των δύο, οπότε βρίσκουμε και τα δύο είδη.

\tmima{Ασκήσεις}

\zoni{Το κατάλαβα}{ποιο σχήμα;}

\ex Σωστό ή λάθος; α) Κάθε τετράγωνο είναι ρόμβος. β) Κάθε ρόμβος είναι τετράγωνο. γ) Οι διαγώνιοι του ρόμβου είναι κάθετες. δ) Οι διαγώνιοι του ορθογωνίου είναι κάθετες. ε) Το ορθογώνιο έχει τέσσερις άξονες συμμετρίας. στ) Το τετράγωνο έχει τέσσερις άξονες συμμετρίας.

\ex Για κάθε ιδιότητα πες σε ποια από τα τρία σχήματα (ορθογώνιο, ρόμβος, τετράγωνο) ισχύει: α) όλες οι γωνίες ορθές β) όλες οι πλευρές ίσες γ) οι διαγώνιοι ίσες δ) οι διαγώνιοι κάθετες ε) οι διαγώνιοι διχοτομούν τις γωνίες στ) τέσσερις άξονες συμμετρίας.

\ex Σχεδίασε: α) ένα ορθογώνιο $5$ cm επί $3$ cm, β) έναν ρόμβο με πλευρά $3$ cm και μία γωνία $60^\circ$, γ) ένα τετράγωνο με πλευρά $4$ cm. Φέρε τις διαγωνίους τους και πες σε τι τρίγωνα χωρίζονται.

\ex Ένας ρόμβος έχει $\hat{\mathrm{Α}}=70^\circ$. Βρες τις άλλες γωνίες του.

\ex Βρες: α) την περίμετρο ορθογωνίου με πλευρές $12$ cm και $7$ cm, β) την πλευρά τετραγώνου με περίμετρο $36$ cm, γ) την πλευρά ρόμβου με περίμετρο $30$ cm.

\ex Στο ορθογώνιο $\mathrm{ΑΒΓΔ}$ η διαγώνιος $\mathrm{ΑΓ}$ είναι $14$ cm. Βρες τη $\mathrm{ΒΔ}$ και τα $\mathrm{ΟΑ}$, $\mathrm{ΟΒ}$, $\mathrm{ΟΓ}$, $\mathrm{ΟΔ}$.

\ex Στον ρόμβο $\mathrm{ΑΒΓΔ}$ είναι $\hat{\mathrm{Α}}=50^\circ$. Βρες τις γωνίες $\widehat{\mathrm{ΟΑΒ}}$ και $\widehat{\mathrm{ΟΒΑ}}$.

\zoni{Το χρησιμοποιώ}{αυλές, παράθυρα, χαρταετοί}

\ex \textbf{Η αυλή.} Η αυλή της θείας Μαρίας είναι τετράγωνη με πλευρά $4{,}5$ m. Πόσα μέτρα περιφράκτη χρειάζεται και πόσους άξονες συμμετρίας έχει το σχήμα της;

\ex \textbf{Το παράθυρο.} Ο παππούς Γιάννης ελέγχει αν ένα παραλληλόγραμμο παράθυρο είναι ορθογώνιο μετρώντας τις διαγωνίους του: $130$ cm και $128$ cm. Είναι ορθογώνιο; Γιατί;

\ex \textbf{Ο χαρταετός.} Ο Νίκος φτιάχνει χαρταετό σε σχήμα ρόμβου με πλευρά $40$ cm. Οι ξύλινες ράβδοι πάνω στις διαγωνίους είναι $64$ cm και $48$ cm. Βρες την περίμετρο του χαρταετού και την περίμετρο ενός από τα τέσσερα τρίγωνα που σχηματίζονται.

\ex \textbf{Τα πλακάκια.} Η Φανή τοποθετεί $5$ τετράγωνα πλακάκια πλευράς $25$ cm σε μία σειρά. Τι σχήμα σχηματίζουν και ποια είναι η περίμετρός του;

\ex \textbf{Ο ισχυρισμός του Στέλιου.} Ο Στέλιος λέει: «Οι διαγώνιοι του ρόμβου είναι κάθετες, άρα και οι διαγώνιοι του ορθογωνίου». Έχει δίκιο; Φέρε ένα παράδειγμα με ορθογώνιο $5$ cm επί $3$ cm.

\zoni{Το σκέφτομαι}{γιατί και πώς}

\ex Εξήγησε με αναδίπλωση γιατί οι διαγώνιοι του ρόμβου τέμνονται κάθετα.

\ex Στο τετράγωνο $\mathrm{ΑΒΓΔ}$ εξήγησε γιατί η διαγώνιος $\mathrm{ΒΔ}$ σχηματίζει γωνία $45^\circ$ με την πλευρά $\mathrm{ΑΒ}$.

\ex Στο ορθογώνιο $\mathrm{ΑΒΓΔ}$ με κέντρο $\mathrm{Ο}$ είναι $\widehat{\mathrm{ΑΒΔ}}=30^\circ$. Βρες τις $\widehat{\mathrm{ΔΒΓ}}$, $\widehat{\mathrm{ΑΟΒ}}$, $\widehat{\mathrm{ΒΟΓ}}$ και πες τι είδους τρίγωνο είναι το $\mathrm{ΟΒΓ}$.

\ex Ένας ρόμβος έχει περίμετρο $52$ cm και μία πλευρά του είναι $3x-2$ cm. Βρες το $x$.

\zoni{Για το σπίτι}{ρίξε μια ματιά}

\ex Κόψε από χαρτί ένα τετράγωνο και ένα ορθογώνιο. Δίπλωσέ τα για να βρεις τους άξονες συμμετρίας τους. Πόσους βρήκες σε κάθε σχήμα;

\ex Βρες στο σπίτι ένα αντικείμενο σχήματος ορθογωνίου, ένα σχήματος τετραγώνου και ένα σχήματος ρόμβου (π.χ. στο σκάκι, σε μοτίβο, σε μπλούζα). Μέτρησε τις διαγωνίους τους και πες ποιες είναι ίσες.
